%% ---------------------------------------------------------------------------
%% intro.tex
%%
%% Introduction
%%
%% $Id: intro.tex 1477 2010-07-28 21:34:43Z palvarado $
%% ---------------------------------------------------------------------------

\chapter{Introducción}
\label{chp:intro}

La Christian-Albrechts-Universitat zu Kiel conocida como CAU o Universidad de Kiel,
referente en investigación tecnológica en Alemania, ha iniciado una investigación para
diseñar una arquitectura de FPGA propia. A diferencia de las soluciones comerciales de
caja negra, este desarrollo busca ofrecer una plataforma competitiva para fines académicos
y experimentales.
El proyecto ha captado el interés de TSMC, el fabricante de semiconductores más grande
del mundo, que ha manifestado su intención de apoyar al centro educativo mediante el
acceso a nodos de fabricación avanzados. Esta sinergia se extiende a Centroamérica gracias
a los lazos históricos de cooperación entre la Universidad de Kiel y el Instituto Tecnológico
de Costa Rica, permitiendo una transferencia de conocimiento de alto nivel en el área de
sistemas digitales complejos.
La investigación se centra en el concepto de la caracterización de una FPGA embebida
en lugar de la fabricación de un IC desde cero. Su desarrollo consta del diseño de una
eFPGA en conjunto de un sistema para la inyección de flujo de bits para su correcta
programación. La esencia del proyecto es hardware reconfigurable dentro de otro para
su caracterización y desarrollo, la implementación de un filtro digital parametrizable y multiplicador de matrices se
requieren para términos del comportamiento de la eFPGA, entre las métricas a estudiar
se encuentra el consumo de potencia, las rutas críticas y la eficiencia de las tablas de
consulta.

\section{Planteamiento del problema}

\subsection{Estado Actual}
 Aunque existen arquitecturas de FPGA de desarrollo académico, los mismos no cuentan
con la capacidad de la flexibilidad de una FPGA comercial debido a características como
la cantidad de tablas de consulta, potencia y velocidad de procesamiento, por lo que su 
desarrollo se ha quedado estancado en su propia academia. Las FPGAs comerciales
requieren licencias para su uso además del elevado precio que significa la compra de
suficientes para impartir un curso. Además, existe un interés de TSMC para su desarrollo,
sin embargo, aún no se cuenta con un modelo caracterizado ni validado.

\subsection{Definición del Problema}

La dependencia de arquitecturas de hardware cerrado limita la capacidad de los centros
de investigación para innovar en el diseño de semiconductores. En la colaboración CAU
TEC, el problema radica en la falta de una metodología validada para caracterizar el
rendimiento, área y consumo de una eFPGA. Sin una implementación previa en una
FPGA física que actúe como entorno de prueba para aplicaciones como filtros digitales,
no es posible asegurar que el diseño final sea competitivo para su futura fabricación con
TSMC.

\subsection{Justificación formal}

La Universidad de Kiel ha identificado la necesidad de desarrollar una arquitectura 
de FPGA propia para fines académicos y experimentales, buscando superar las barreras 
económicas y de licenciamiento que imponen las soluciones comerciales de hardware cerrado. 
Sin embargo, para que este desarrollo sea viable y competitivo frente a una futura 
fabricación en silicio con nodos avanzados, es necesario validar previamente su diseño.

Existe una carencia de documentación y metodologías validadas sobre 
los compromisos entre parámetros estructurales clave en arquitecturas de hardware abierto. 
Por lo tanto, la realización de este proyecto se justifica en la necesidad técnica de proveer 
una prueba de concepto y una metodología de emulación física que permita caracterizar el 
comportamiento de la unidad de lógica reconfigurable antes de su fabricación. Al determinar 
los parámetros de funcionamiento como la eficiencia de área por conglomerado de tablas de 
búsqueda, los retardos de temporización y el consumo de potencia, se mitigarán los riesgos 
financieros y de diseño asociados a la impresión en silicio, lo cual dará mayor dirección al diseño 
con respecto a los parámetros deseados.

\section{Estructura del Documento}

En el capítulo 2 se presenta un informe de la investigación del 
estado del arte actual del problema, así como una descripción de las partes y 
características del proyecto. Se profundiza en la arquitectura básica de una FPGA 
y en el método para cargarle el bitstream.

En el capítulo 3 se describe el proceso completo de diseño del proyecto, 
desde la definición de los requisitos iniciales hasta la obtención del sistema final. 
Se justifican las decisiones tomadas en cada etapa, así como las alternativas que se 
valoraron y los motivos por los que fueron descartadas. Además, se documenta el flujo 
de trabajo seguido, con las herramientas empleadas y los pasos realizados, de modo que 
el desarrollo pueda ser comprendido y reproducido.

%\setlength{\parskip}{2.5mm}
%
%En el \textbf{capítulo~\ref{ch:solucion}} se presenta las posible solución al problema 
%donde analizamos las etapas a tomar en cuenta para la validación de una arquitectura de 
%lógica reconfigurable, se describe además cada parte del proyecto y los resultados que 
%se esperan.

%\setlength{\parskip}{2.5mm}

%En el \textbf{capítulo~\ref{ch:HyO}}  se expone la hipótesis planteada, además, se 
%presenta el objetivo general para este trabajo, así como, los objetivos específicos 
%del mismo. Se muestran metas por objetivo propuesto. Por último, 
%se resaltan los alcances y limitaciones que presenta este trabajo.

%\setlength{\parskip}{2.5mm}

%En el \textbf{capítulo~\ref{ch:Experimento}} se presenta el posible experimento 
%preliminar que detalla la estrategia del proyecto y el enfoque necesario prueba el 
%sistema y el flujo de herramientas a través de un experimento bien diseñado y riguroso.

%\setlength{\parskip}{2.5mm}

%En el \textbf{capítulo~\ref{ch:Crono}}  se presentan las tareas necesarias a llevar a 
%cabo a lo largo de este trabajo, así como el diagrama de Gantt para poder completar con las 
%tareas propuestas.

\setlength{\parskip}{2.5mm}

\section{Hipótesis}
     ``La variación de la cantidad de bloques DSP y de la altura del fabric en una 
     eFPGA de hardware abierto, manteniendo constante la cantidad de azulejos LUT4AB, 
     produce diferencias medibles en la utilización de recursos, los parámetros de 
     temporización y la potencia estimada al implementar un filtro FIR y un multiplicador 
     de matrices."

\section{Objetivos}

\subsection{General}
            Evaluar el impacto de la variación de parámetros estructurales en el 
            rendimiento, área y potencia de una arquitectura abierta de lógica 
            reconfigurable a nivel de simulación y sintesis utilizando una 
            FPGA comercial como plataforma de emulación física.
\subsection{Específicos}
	\begin{itemize}
                \item Seleccionar una arquitectura de lógica reconfigurable de hardware 
                abierto y su flujo de herramientas asociado, justificando la elección 
                basándose en su nivel de parametrización, soporte de verificación funcional 
                e historial de implementación física en otros proyectos.

                \item Configurar la arquitectura seleccionada en una tarjeta 
                FPGA, basándose en reportes de síntesis, ruteo y ejecución 
                funcional, mediante el uso de al menos dos circuitos de prueba en dicha 
                lógica programable.
                             
                \item Caracterizar el impacto de la variación de los parámetros 
                estructurales disponibles en la arquitectura de lógica programable sobre 
                la utilización de recursos, el consumo de potencia y el rendimiento a nivel 
                de temporización.
    \end{itemize}


    \begin{table}[htbp]
\small
\begin{tabularx}{\textwidth}{|p{1.8cm}|X|X|X|}
\hline
\textbf{Etapas} & \textbf{Problemas} & \textbf{Hipótesis} & \textbf{Objetivo Específico} \\ \hline

\textbf{Selección} & 
Limitada capacidad de innovación debido a la dependencia de arquitecturas cerradas y la falta de criterios claros para elegir un flujo abierto parametrizable. & 
La selección de una arquitectura abierta con alto nivel de parametrización e historial físico garantiza un flujo de verificación funcional confiable. & 
Seleccionar una arquitectura de lógica reconfigurable de hardware abierto y su flujo de herramientas asociado, justificando la elección en su parametrización e historial. \\ \hline

\textbf{Emulación} & 
Incertidumbre sobre la competitividad del diseño final ante la falta de una plataforma de emulación física previa para circuitos de prueba. & 
La instanciación de la eFPGA en una tarjeta física comercial permite comprobar su viabilidad funcional ejecutando circuitos de prueba como filtros digitales o multiplicadores. & 
Configurar e instanciar la arquitectura seleccionada en una tarjeta FPGA comercial, caracterizarla funcionalmente e implementar al menos dos circuitos de prueba. \\ \hline

\textbf{Métricas} & 
Carencia de documentación y metodologías validadas sobre cómo afectan los parámetros estructurales los compromisos de área, retardo y potencia. & 
La variación de la cantidad de bloques DSP y de la altura del fabric produce diferencias medibles en la utilización de recursos, la temporización y la potencia estimada. & 
Caracterizar el impacto de la variación de los parámetros estructurales disponibles sobre la utilización de recursos, el consumo de potencia y el rendimiento temporal. \\ \hline
\end{tabularx}
\caption{Matriz de trazabilidad del proyecto.}
\label{tab:matriz_trazabilidad}
\end{table}
 
\section{Lista de Tareas por Objetivo}

\textbf{\normalsize{Objetivo Específico 1}}

Este objetivo se centra en la investigación de la arquitectura de 
lógica reconfigurable y el flujo de herramientas para su implementación, 
considerando la el soporte para su parametrización, verificación 
funcional e historial de implementación física.

\begin{enumerate}

    \item Investigar y recopilar información sobre las distintas 
    arquitecturas de lógica reconfigurable de hardware abierto 
    disponibles.
     
    \item Evaluar la flexibilidad y el nivel de parametrización 
    estructural que ofrece cada arquitectura preseleccionada.
    
    \item Analizar la madurez del flujo de herramientas asociado, 
    comprobando la compatibilidad y el soporte para la verificación 
    funcional.

    \item Revisar el historial de implementación física en silicio 
    de las opciones evaluadas para constatar su viabilidad y 
    confiabilidad en proyectos previos. 

    \item Seleccionar la arquitectura definitiva y justificar técnicamente 
    la elección, basada en el cumplimiento de los criterios previamente 
    analizados.

\end{enumerate}

Las tareas descritas justifican el uso de la arquitectura en las próximas 
etapas y su implementación en la FPGA.

\textbf{\normalsize{Objetivo Específico 2}}

El segundo objetivo se centra en la configuración de la arquitectura 
seleccionada mediante JTAG y la caracterización basada en reportes 
desplegados por las herramientas seleccionadas en la etapa anterior, 
además de la implementación de dos circuitos de prueba.

\begin{enumerate}

    \item Ajustar los módulos de la arquitectura seleccionada para asegurar 
    su compatibilidad con la herramienta de síntesis de la FPGA objetivo.
    
    \item Integrar la arquitectura en el proyecto principal de la herramienta de desarrollo.
    
    \item Correr el proceso de síntesis y extraer los reportes de la cantidad de LUTs, 
    flip-flops y bloques de memoria requeridos.

    \item Ejecutar la implementación en la FPGA y compilar los reportes de tiempos 
     de setup y hold, consumo de potencia estimado y congestión de ruteo.

    \item Generar el bitstream necesario para configurar la FPGA física.
    
    \item Seleccionar dos circuitos que estresen diferentes aspectos de la lógica programable.
    
    \item Codificar ambos circuitos independientemente.
    
    \item Escribir los testbenches correspondientes y simular el comportamiento 
    de los circuitos de prueba para asegurar su corrección funcional antes de la 
    implementación física.

    \item Compilar y cargar el primer y segundo circuito de prueba sobre la 
    arquitectura de lógica programable ya instanciada.

    \item Estimular las entradas en el hardware real, capturar las salidas y contrastar la 
    ejecución funcional contra los resultados obtenidos en las simulaciones

\end{enumerate}

Estas tareas permiten validar la arquitectura seleccionada y dan información benéfica para 
el desarrollo de la caracterización de la arquitectura seleccionada.

\textbf{\normalsize{Objetivo Específico 3}}

El tercer objetivo da base a la caracterizacion de la arquitecura de lógica programable 
seleccionada según la utilización de recursos.

\begin{enumerate}

    \item Identificar cuáles variables de la arquitectura de la arquitectura lógica 
    programable serán modificables.

    \item Definir un conjunto discreto de valores para cada parámetro seleccionado y 
    limitar las combinaciones posibles a un número manejable y representativo.
    
    \item Recopilar la cantidad de bloques lógicos básicos utilizados y la estimación de 
    área equivalente de silicio.

    \item Extraer del Análisis de Temporización Estática y el retardo de la ruta crítica.

    \item Compilar las estimaciones de consumo de potencia dinámica y estática reportadas 
    tras el ruteo.

    \item Procesar los datos consolidados para graficar el comportamiento de las métricas 
    frente a la variación de cada parámetro.

    \item Analizar los puntos donde la mejora de una métrica perjudica a otra (Trade-off).
    
    \item Concluir cuáles son las configuraciones de parámetros estructurales que ofrecen 
    el mejor balance técnico para los casos de uso proyectados.

\end{enumerate}


\section{Alcances y Limitaciones del Proyecto}

En la sección de limitaciones se exponen los desafíos y restricciones propios de la 
naturaleza del proyecto. Si bien estas condiciones pueden representar obstáculos, 
al mismo tiempo delimitan el alcance del trabajo y establecen un marco de referencia 
realista para su ejecución. A continuación, se detallan las principales limitaciones 
identificadas:

\begin{itemize}

    \item La ejecución de los experimentos se llevará a cabo con la FPGA 
    disponible durante la evaluación del proyecto. Esta limitación en los recursos 
    físicos para la realización de pruebas podría restringir tanto la diversidad de 
    los escenarios evaluados como la naturaleza de los datos susceptibles de recolección.

    \item Al instanciar una arquitectura de lógica programable sobre una FPGA comercial, 
    el área, la potencia y el retardo de temporización estarán limitados por la arquitectura 
    de la FPGA anfitriona.

    \item La configuración óptima identificada estará restringida a las características 
    de los circuitos seleccionados y podría no ser generalizable a todas las aplicaciones 
    digitales.

    \item No se realizará una exploración exhaustiva de todo el espacio de diseño, por lo 
    que la configuración óptima que maximiza el rendimiento corresponderá únicamente a un 
    óptimo dentro del subconjunto explorado.

    \item Las métricas de consumo de potencia se extraerán a partir de reportes generados 
    por herramientas de síntesis y ruteo, no mediante mediciones físicas directas sobre 
    la arquitectura implementada por lo tanto se introduce un margen de error inherente a 
    los modelos matemáticos y estimaciones.



\end{itemize}

El alcance de este trabajo de graduación se limita a la etapa de diseño conceptual, 
caracterización y validación metodológica de tres arquitecturas de lógica reconfigurable 
propuestas. El objetivo principal no es la entrega de un dispositivo físico final, sino 
la obtención de un marco de evaluación cuantitativo que permita determinar la viabilidad 
técnica y los trade-offs de rendimiento de diversas configuraciones de lógica reconfigurable.

En este sentido, la investigación se limita a la fase de simulación y síntesis lógica mediante 
el uso de herramientas de código abierto como Yosys y FABulous, sin embargo, para la programación
de la FPGA, se utilizará Quartus® al igual que para la obtencion del consumo de potencia. El proyecto culminará con la 
entrega de los resultados experimentales de área, retardo y consumo de potencia, los cuales 
son derivados de la implementación de los benchmarks seleccionados. Por tanto, este trabajo 
constituye una prueba de concepto fundamental que sienta las bases técnicas y metodológicas 
para que en etapas posteriores, los resultados sirvan como guía y referencia para el diseño 
físico definitivo y la eventual fabricación en silicio de la arquitectura analizada. De esta 
manera, el alcance del presente estudio se centra exclusivamente en la caracterización teórica 
y la optimización de los parámetros de diseño, proporcionando la dirección necesaria para el 
desarrollo de futuras implementaciones en nodos avanzados de fabricación.

\section{Solución Propuesta}


Este proyecto propone el desarrollo de una metodología de emulación física y caracterización para 
superar las limitaciones del hardware cerrado y disminuir los riesgos de la fabricación en 
silicio sin el previo modelado del sistema necesario. 
Para esto, se proponen cuatro etapas para la solución.

\subsection{Primera Etapa}

La primera etapa de la solución propuesta para el proyecto se enfoca en la selección de 
una arquitectura de lógica reconfigurable de código abierto parametrizable, y su 
implementación mediante un banco de pruebas para comprobar su funcionamiento correcto. De 
esta manera se reduce la dependencia del flujo de herramientas de un fabricante específico 
y se asegura la portabilidad del proyecto a diferentes plataformas, al igual que se promueve 
la continuidad del proyecto en el caso de querer profundizar en el futuro.

\subsection{Segunda Etapa}

La segunda parte de la solución propuesta para el proyecto se centra en la instanciación de 
la arquitectura de lógica reconfigurable seleccionada en la primera etapa sobre una FPGA 
comercial en Quartus®. Mediante esta etapa se busca un medio para extraer en próximas etapas las métricas 
de rendimiento y uso de recursos como el uso de los recursos, el consumo de potencia y los 
retardos de las rutas críticas.

\subsection{Tercera Etapa}

La tercera fase de la solución propuesta para el proyecto busca implementar dos circuitos 
digitales que estresen la arquitectura de lógica reconfigurable seleccionada en la primera 
etapa de maneras diferentes. Para esto se propone la implementación de un filtro digital 
de orden bajo y un multiplicador de matrices. De esta manera se busca obtener métricas de 
rendimiento para su comparación y se garantiza la funcionalidad del flujo de síntesis y 
ruteo de la arquitectura de lógica reconfigurable seleccionada.

\subsection{Cuarta Etapa}

La cuarta etapa de la solución propuesta para el proyecto tiene como objetivo la consolidación 
de los datos obtenidos en la tercera etapa observando cómo las métricas de rendimiento se 
compensan entre sí, para determinar si la arquitectura de lógica reconfigurable seleccionada 
es adecuada para su uso en aplicaciones de alto rendimiento o si es necesario realizar 
modificaciones según la aplicación en la que se utilice. Esta etapa se enfoca en la 
interpretación de los datos y cómo su análisis puede guiar el diseño de futuras arquitecturas 
de lógica reconfigurable para su impresión en silicio.

De esta manera se establece que para validar la viabilidad y el rendimiento de la arquitectura 
eFPGA propuesta, se definió el marco experimental que permite la caracterización del diseño bajo 
condiciones controladas. Se fundamenta en la evaluación comparativa de tres arquitecturas de 
lógica reconfigurable distintas, las cuales varían en sus parámetros estructurales, como las dimensiones de la arquitectura y la cantidad de bloques DSP. Las configuraciones serán 
sometidas a pruebas mediante la implementación de dos circuitos de referencia que actuarán 
como criterios de evaluación, un multiplicador de matrices, seleccionado por su alta 
exigencia de recursos aritméticos y capacidad para estresar la estructura de interconexión, 
y un filtro digital, empleado para evaluar el comportamiento de la arquitectura en tareas de 
procesamiento de señales y latencia de procesamiento secuencial.

La cuantificación del impacto de cada configuración se realizará mediante figuras de mérito 
que permitan un análisis objetivo. Se medirá el área mediante el conteo de recursos lógicos equivalentes tras la síntesis, en segundo lugar, se 
determinará parametros de temporización para establecer la frecuencia 
máxima de operación del sistema, y se calculará el consumo de potencia en miliwatts, 
incluyendo la estimación de potencia.

El éxito de la solución propuesta se determinará mediante la implementación exitosa de los 
benchmarks sobre cada una de las tres arquitecturas. El criterio de evaluación principal 
consiste en el análisis comparativo de las métricas obtenidas para cada combinación de 
arquitectura y diseño de prueba. De esta manera, se busca determinar cómo varían el uso del 
recursos, la temporización y el consumo energético en función de la estructura de lógica reconfigurable, 
estableciendo así un balance óptimo que demuestre la funcionalidad y eficiencia del diseño 
antes de su proyección a fabricación en silicio.
%%% Local Variables: 
%%% mode: latex
%%% TeX-master: "main"
%%% End: 
